\section{Data Consumption, Quantiles, and Auditing}
\label{sec:design}
Before execution, the allocator uses task attributes and histories from
$\mathcal{P}(t)$ to issue $\hat y_t$ (Section~\ref{sec:problem}). Auditing
collects consumption during execution; after completion, recorded
measurements enter the histories used for subsequent requests.

\subsection{Workflow integration, measurement, and history}
\label{sec:history}
A workflow engine such as Nextflow can call an allocator before submission.
A tracer follows the task launcher and its child processes; completion
callbacks link memory and consumption measurements to the task identifier.
This integration leaves analysis programs unchanged.

On Linux, eBPF handlers filter events by the tracked process tree.
Explicit-read accounting associates file descriptors with observed file opens
and sums bytes returned by \texttt{read} and \texttt{pread64}.
A successful \texttt{mmap} call creates a virtual mapping, not a measurement
of bytes accessed. File-backed page faults, observed for example at
\texttt{filemap\_fault}, provide a page-based access estimate. Combined with
explicit-read bytes, this signal describes access activity, not unique
bytes or peak memory.

A compatibility scope is a group of earlier executions relevant to a
pending task. Scopes may match task instances, inputs, static features,
processes, workflows, or execution systems, preserving software and
resource context. For example, a repeated alignment task can use matching
input history; a new input may require broader process history. Minimum
observation counts govern fallback. Peak-memory summaries use completed
tasks; consumption summaries use their audited subset.

\subsection{Consumption features and prelaunch prediction}
\label{sec:consumption}
For task $t$, let $\mathbf{s}_t$ contain static attributes,
$\mathbf{h}_{m,t}$ summarize earlier peak memory, and $\mathbf{h}_{c,t}$
summarize earlier audited consumption, as defined in Section~\ref{sec:problem}.
An auxiliary estimator $g$ predicts consumption
$\hat c_t=g(\mathbf{s}_t,\mathbf{h}_{m,t})$, yielding memory-model input
\begin{equation}
\mathbf{z}_t = [\mathbf{s}_t,\mathbf{h}_{m,t},\mathbf{h}_{c,t},\hat c_t].
\label{eq:general_features}
\end{equation}
The estimator learns from completed audited tasks. The pending task's
measured consumption is never a prelaunch input. For offline training,
grouped out-of-fold estimates predict each group of related tasks using
an estimator trained on other groups, avoiding use of that group's
consumption targets to construct its predictions.

\label{sec:quantiles}
A quantile model estimates $\widehat Q_\alpha(\mathbf{z}_t)$, the memory
percentile at level $\alpha$. The median describes central demand; an
upper percentile targets a smaller exceedance probability. For example,
$\alpha=0.95$ targets a threshold exceeded by approximately 5\% of comparable
tasks, not guaranteed 95\% coverage on new inputs. Models may specialize
by process, with broader fallbacks. Holding the prediction configuration
fixed when adding either consumption feature or both isolates their contributions.

\subsection{Calibration and budgeted selective feedback}
\label{sec:calibration}
Calibration adjusts a candidate quantile using completed-task prediction errors.
For example, fitting a positive multiplier $k_t$ from earlier
observed-to-predicted memory ratios gives a request
\begin{equation}
\hat y_t=\left\lceil k_t\widehat Q_\alpha(\mathbf{z}_t)\right\rceil_{\mathrm{MiB}}.
\label{eq:general_request}
\end{equation}
Here $\hat y_t$ is the initial request, rounded to whole mebibytes.
Prediction and calibration together implement $f$ in Section~\ref{sec:problem}.
Development data determine the quantile and correction policy; only
outcomes available before submission inform $k_t$. Compatible prediction-error
histories and broader fallbacks support calibration, whose achieved
coverage can change when inputs shift.

\label{sec:gates}
A gating process applies selection rules to prelaunch predictions and
the available historical observations. It selects tasks whose consumption features are expected
to improve memory predictions for similar subsequent tasks. Over $N$ tasks,
binary audit decisions $q_t\in\{0,1\}$ must satisfy
$\sum_{t=1}^{N}q_t\leq\lfloor BN\rfloor$, where $B$ is the allowed audit
fraction and $q_t=1$ selects a task.
The budget limits measurement volume; gates determine which tasks receive
it. Completed peak-memory observations update history, but consumption
enters only when audited. Feedback benefits later requests, not the
request already issued. Equal-budget comparisons and measurement costs
are therefore needed to assess selection quality.

\subsection{CAMP implementation and relationship to WORKS}
\label{sec:camp-implementation}
Our WORKS paper implemented these mechanisms in CAMP~\cite{camp2026works}:
task-attributed eBPF measurements, compatibility scopes, a consumption
estimator, and budgeted selection. LightGBM ensembles estimate the 50th,
90th, 95th, 99th, and 99.5th memory percentiles. The process-tail configuration
uses a global median model, workflow--process upper-tail models, and
global fallbacks. Residual calibration blends compatible and broader
histories, permitting only upward corrections. These mechanisms and the
published component results are prior work. Here we expand retained
component and auditing-cost analyses~\cite{campArtifact} and add distinct
live transfer experiments on unseen inputs.
